% Generated draft: requires mathematical and source review.
\documentclass[11pt]{article}
\usepackage[margin=25mm]{geometry}
\usepackage{fontspec}
\setmainfont{Latin Modern Roman}
\usepackage{amsmath,amssymb,mathtools,unicode-math}
\setmathfont{Latin Modern Math}
\usepackage{enumitem,needspace,xurl,hyperref}
\setlength{\parindent}{0pt}
\setlength{\parskip}{5pt}
\setlength{\emergencystretch}{3em}
\newcommand{\N}{\mathbb{N}}
\newcommand{\Z}{\mathbb{Z}}
\newcommand{\Q}{\mathbb{Q}}
\newcommand{\R}{\mathbb{R}}
\newcommand{\C}{\mathbb{C}}
\newcommand{\F}{\mathbb{F}}
\newcommand{\A}{\mathbb{A}}
\newcommand{\PP}{\mathbb{P}}
\newcommand{\E}{\mathbb{E}}
\newcommand{\eps}{\varepsilon}
\newcommand{\dd}{\,\mathrm{d}}
\begin{document}
\section*{UnsolvedMath ID: 30005271}
\textbf{Source problem code:} catalogue entry
\subsection{Definitions and mathematical statement}
For a positive integer \(n\), define the Möbius function by
\[
\mu(n)=\begin{cases}
1,&n=1,\\
(-1)^k,&n\text{ is a product of }k\text{ distinct primes},\\
0,&n\text{ is divisible by the square of a prime}.
\end{cases}
\]
For real \(x>e\) and \(y>0\), put
\[
M(x,y)=\sum_{\substack{n\in\Z\\x<n\le x+y}}\mu(n).
\]
The supplied problem asks for cancellation in \(M(x,y)\) at the square-root interval scale or below. It singles out the range
\[
\frac{\sqrt{x}}{(\log\log x)^{1/4}}\ll y\le x^{1-\epsilon},
\]
where \(\log\) denotes the natural logarithm. The lower-bound notation means that \(\sqrt{x}/(\log\log x)^{1/4}\le K y\) for a positive constant \(K\) independent of \(x\) and \(y\); its permitted dependence on other parameters is not specified.

\textbf{Unresolved formulation.} The record does not define the required strength of cancellation, the quantification of \(\epsilon\), or the uniformity in \(y\). A possible precise interpretation, presented for review rather than as an established transcription, is: for every fixed \(0<\epsilon<1/2\) and every fixed \(c>0\),
\[
\lim_{x\to\infty}\ \sup_{\displaystyle c\sqrt{x}/(\log\log x)^{1/4}\le y\le x^{1-\epsilon}}
\frac{|M(x,y)|}{y}=0.
\]
This interprets cancellation as a uniform \(o(y)\) bound. The supplied evidence does not determine whether this interpretation, a quantitative saving, or an assertion restricted to most intervals is intended. The displayed restriction on \(\epsilon\) makes this illustrative range nonempty for all sufficiently large \(x\); it is not an explicit source hypothesis.

The record attributes the concrete range to a proposed fractional-moment bound but supplies neither that bound nor its hypotheses. Its original quotation describes an unspecified conjecture as plausibly attackable, at least under standard conjectures such as the Riemann Hypothesis, and ends before the claimed consequence of Perron's formula. Consequently, neither an unconditional target nor a precise conditional implication can be recovered from that quotation.
\subsection{Short English statement}
The problem seeks cancellation in sums of the Möbius function over short intervals, including lengths near \(\sqrt{x}/(\log\log x)^{1/4}\). The supplied record leaves the strength, uniformity, and hypotheses of the desired estimate unresolved. The displayed uniform \(o(y)\) statement is an explicitly provisional interpretation requiring human mathematical review.
\subsection{Sources}
\begin{itemize}
\item Local UnsolvedMath registry, \texttt{data/unsolvedmath/problems.json}, supplied record with permanent ID \(30005271\), exact source problem code \texttt{catalogue entry}, title ``Square-Root-Scale Cancellation in Short Möbius Sums.'' The record gives proposed year 2022 and supplies no named proposer.
\item Supplied report citation: \emph{Analytic Number Theory} (2023), Oberwolfach Reports, pp.~2895--2960, DOI \texttt{10.4171/owr/2022/50}. \url{https://doi.org/10.4171/owr/2022/50}. The background describes Oberwolfach Reports as workshop reports documenting mathematical discussions and open problems. The supplied citation year is 2023, while the DOI contains 2022; this discrepancy is retained without resolution.
\item Supplied literature reference: Guth and Maynard. \url{https://doi.org/10.4007/annals.2026.203.2.6}.
\item Supplied literature reference: Oberwolfach Report 50. \url{https://ems.press/content/serial-article-files/46986}.
\item Literature assessment date: \texttt{2026-08-22}. This is supplied dated curation, not independently verified. The record labels the problem partially solved and reports that Guth and Maynard improve the relevant prime-interval exponent to \(17/30\), up to lower-order terms, without reaching the square-root barrier. The record itself requests independent verification of time-sensitive status and cited claims. No precise implication for the Möbius-sum target is supplied.
\item The registry labels the statement \texttt{corrected\_verified} and says it was reconstructed and checked against the official report. These are supplied provenance claims; this draft does not independently verify them or the mathematical formulation.
\end{itemize}
\end{document}
